% vim: set ft=tex tabstop=4 shiftwidth=4 noexpandtab:

\input{../../include/latex/tikz/header.tex}

% opening %{{{1

\begin{document}
\begin{tikzpicture}[scale=1.0]

% parameters %{{{1

	\tikzmath{
		\radius = 5.0 ;
		\rotation = 90 ;
		\numsidesfloat = 6.0 ;
		\numsides = int(\numsidesfloat) ;
		\numsidesminus = \numsides - 1 ;
		int \i ;
		for \i in {1,...,\numsides} {
			\ang{\i} = \rotation + 360 / \numsidesfloat * (\i-1) ;
		};
		\xmin = - \radius ;
		\xmax = \radius + 1 ;
		\ymin = \radius/4 ;
		\ymax = \radius + 1 ;
	}

% box %{{{1

	\useasboundingbox (\xmin,\ymin) rectangle (\xmax,\ymax);

% coordinates %{{{1

	\coordinate (O) at (0,0);

	\foreach \i in {1,...,\numsides}
		\coordinate (P_\i) at (\ang{\i}:\radius);

	\foreach \i in {1,...,\numsides}
		\coordinate (I_\i) at (\ang{\i}:{\radius/2});

	\coordinate (I) at (\rotation:{\radius/2});

% intersections %{{{1

	\begin{pgfinterruptboundingbox}
		\path[name path=line13] (P_1) -- (P_3);
		\path[name path=line26] (P_2) -- (P_6);
		\path[name intersections={of=line13 and line26, by=J}];
	\end{pgfinterruptboundingbox}

	\begin{pgfinterruptboundingbox}
		\path[name path=line15] (P_1) -- (P_5);
		\path[name path=line26] (P_2) -- (P_6);
		\path[name intersections={of=line15 and line26, by=K}];
	\end{pgfinterruptboundingbox}

% points, dots, vertices %{{{1

	\fill (O) circle (0.4mm);

	\foreach \i in {1,...,\numsides}
		\fill (P_\i) circle (0.4mm);

	\foreach \i in {1,...,\numsides}
		\fill (I_\i) circle (0.4mm);

	\fill (J) circle (0.4mm);
	\fill (K) circle (0.4mm);

% segments, sides, lines %{{{1

	\draw (P_1) \foreach \i in {2,...,\numsides} { -- (P_\i) } -- cycle;

	\foreach \i in {1,...,\numsides}
		\draw (O) -- (P_\i);

	\draw (P_1) -- (P_3) -- (P_5) -- cycle ;
	\draw (P_2) -- (P_4) -- (P_6) -- cycle ;

% circles %{{{1

	\draw (O) circle (\radius);

% points, dots, vertices labels %{{{1

	\node[anchor=east] at ($ (O) + (180:0.15) $) {$O$};

	\node[above] at (P_1) {$P_1$};
	\node[left] at (P_2) {$P_2$};
	\node[left] at (P_3) {$P_3$};
	\node[below] at (P_4) {$P_4$};
	\node[right] at (P_5) {$P_5$};
	\node[right] at (P_6) {$P_6$};

	\node[below right] at (I) {$I$};
	\node[below right] at (J) {$J$};
	\node[below left] at (K) {$K$};

% segments, sides, lines labels %{{{1

	\node[above left] at ($ (P_1)!0.5!(P_2) $) {$r$};

	\node[above left] at ($ (P_1)!0.5!(J) $) {$a$};

% segments, sides, lines marks %{{{1

	\foreach \i
		[evaluate=\i as \j using { int(mod(\i,\numsides) + 1) }]
		in {1,...,\numsides}
	{
		\tkzMarkSegments[mark=|, size=2pt](P_\i,P_\j);
	}

% circles labels %{{{1

	\node at ($ (O) + (60:\radius + 0.3) $) {$\mathscr{C}$};

% circular arcs labels %{{{2

	%\node at ($ (O) + (45:\radius + 0.3) $) {$\mathscr{C}$};

% angles labels %{{{1

	\pic[draw, ->, "$30^\circ$", angle radius=1.1cm, angle eccentricity=1.3]
	{angle = O--P_1--P_5};
	\pic[draw, ->, "$30^\circ$", angle radius=1.1cm, angle eccentricity=1.3]
	{angle = P_3--P_1--O};

	\pic[draw, ->, "$\alpha$", angle radius=0.8cm, angle eccentricity=1.3]
	{angle = I--J--P_1};

	\pic[draw, ->, "$\beta$", angle radius=0.8cm, angle eccentricity=1.3]
	{angle = P_1--K--I};

	\pic[draw, ->, "$\lambda$", angle radius=0.6cm, angle eccentricity=1.5]
	{angle = P_1--J--P_2};

	\pic[draw, ->, "$\lambda$", angle radius=0.7cm, angle eccentricity=1.5]
	{angle = P_3--J--P_6};

	\pic[draw, ->, "$\mu$", angle radius=0.6cm, angle eccentricity=1.5]
	{angle = P_6--K--P_1};

	\pic[draw, ->, "$\mu$", angle radius=0.7cm, angle eccentricity=1.5]
	{angle = P_2--K--P_5};

% right angles markers %{{{2

	\foreach \i
		[
			evaluate=\i as \j using { int(mod(\i, \numsides) + 1) }
		]
		in {1,...,\numsides}
	{
		\pic [draw, angle radius=7pt, angle eccentricity=1]
		{right angle = O--I_\i--P_\j};
	}

% closing %{{{1

\end{tikzpicture}
\end{document}
